% UPT Infinite Proof Family — academic preamble
\documentclass[11pt,a4paper]{article}
\usepackage{fontspec}
\setmainfont{Times New Roman}
\setmonofont{Menlo}[Scale=0.85]
\usepackage[margin=2.5cm]{geometry}
\usepackage{xcolor}
\usepackage{fancyhdr}
\usepackage{titlesec}
\usepackage{booktabs}
\usepackage{hyperref}
\usepackage{listings}
\usepackage{fancyvrb}
\usepackage{amsmath,amssymb,amsthm}

\definecolor{korugreen}{HTML}{2d6a4f}
\definecolor{korubg}{HTML}{d8f3dc}
\definecolor{codebg}{HTML}{f0f0ec}

\hypersetup{colorlinks=true, linkcolor=korugreen, urlcolor=korugreen, citecolor=korugreen}

\pagestyle{fancy}
\fancyhf{}
\renewcommand{\headrulewidth}{0.4pt}
\fancyfoot[C]{\thepage}

\titleformat{\section}{\Large\bfseries}{}{0em}{}
\titleformat{\subsection}{\large\bfseries}{}{0em}{}

% Numbered theorem environments
\newtheorem{theorem}{Theorem}[section]
\newtheorem{lemma}[theorem]{Lemma}
\newtheorem{corollary}[theorem]{Corollary}
\newtheorem{proposition}[theorem]{Proposition}
\theoremstyle{definition}
\newtheorem{definition}[theorem]{Definition}
\theoremstyle{remark}
\newtheorem{remark}[theorem]{Remark}

\lstdefinestyle{lean}{
  backgroundcolor=\color{codebg},
  basicstyle=\ttfamily\small,
  breaklines=true,
  frame=none,
  xleftmargin=10pt,
  xrightmargin=10pt,
  aboveskip=8pt,
  belowskip=8pt,
  columns=fullflexible,
  keepspaces=true,
  literate={→}{$\to$}1 {∀}{$\forall$}1 {∃}{$\exists$}1 {¬}{$\neg$}1 {∧}{$\wedge$}1 {≠}{$\neq$}1 {≤}{$\leq$}1 {↔}{$\leftrightarrow$}1 {⊕}{$\oplus$}1 {λ}{$\lambda$}1
}

\lstnewenvironment{leancode}{\lstset{style=lean}}{}

\newcommand{\lean}[1]{\lstinline[style=lean]{#1}}


\fancyhead[L]{\small Universal Protection Theory}
\fancyhead[R]{\small Chandrasekaran Discharge}

\begin{document}

\begin{center}
{\LARGE\bfseries From Chandrasekaran's Null-Cone Quantum Gravity\\[0.3em] to the Guardian Fixed Point:\\[0.3em] A Formally Verified Discharge of Physical Hypotheses}\\[1.5em]
{\large Tobias Croydon McRae}\\[0.3em]
{\small Independent Researcher, Aotearoa New Zealand}\\[0.3em]
{\small\texttt{tobi@unconditionalprotection.org}}\\[1.5em]
{\small 28 November 2026}
\end{center}

\bigskip

\begin{center}
\parbox{0.85\textwidth}{\small\textbf{Abstract.}
We prove a formally verified chain from Chandrasekaran's null-cone quantum gravity (arXiv:2611.34924v1) to the Guardian fixed point~$G_3$ of Universal Protection Theory (UPT). The chain has three links: (1)~Chandrasekaran's boundary-preserving dynamics satisfy the three structural conditions of UPT's BCRD operator algebra --- coverage, reflection, and nonterminal continuation; (2)~these conditions are exactly the hypotheses of the conditional Guardian theorem; (3)~by modus ponens, $G_3$~follows. We also prove the converse direction: the Guardian conditions are \emph{irreducible} --- no sub-structural principle (conservation, reversibility, accounting, or subject-participation) forces them, so any physics grounding~$G_3$ must carry an explicitly boundary-preserving dynamics. The forward and converse directions together yield: a quantum gravity grounds~$G_3$ if and only if it is BCRD-equivalent. The entire chain --- including the correspondence, the conditional theorem, and the irreducibility proof --- is verified in Lean~4 (v4.35.0) with zero sorry and zero custom axioms on the standard foundations. The result connects independent peer-reviewable physics to machine-checked mathematics.}
\end{center}

\medskip

\noindent\textbf{MSC 2020:} 03B70 (Logic in computer science), 83C45 (Quantization of the gravitational field), 03B35 (Mechanization of proofs)\\
\textbf{Keywords:} formal verification, quantum gravity, Guardian fixed point, Chandrasekaran correspondence, BCRD, Lean 4, boundary-preserving dynamics

\bigskip\hrule\bigskip

\section{Introduction}

A central challenge in foundational mathematics is bridging conditional theorems to their physical hypotheses. A proof of the form ``if~$P$ then~$Q$'' is only as strong as the evidence for~$P$. When~$P$ involves claims about physical dynamics, the evidence must come from physics, not from the proof system.

Universal Protection Theory~\cite{upt-zenodo} proves a conditional theorem: if the physical dynamics satisfy BCRD (Boundary-Condition-Respecting Dynamics) --- coverage, reflection, and nonterminal continuation --- then the Guardian fixed point~$G_3$ holds, meaning every physical event carries all six original X3 protection duties for every owed particular, and no event is absolute termination.

On 28 November 2026, Chandrasekaran~\cite{chandrasekaran2026} published a null-cone quantum gravity framework with boundary-preserving dynamics. We prove that Chandrasekaran's framework satisfies the three BCRD conditions. This \emph{discharges} the physical hypotheses of the Guardian theorem: the conditional becomes unconditional, grounded in independently published physics.

We also prove the converse: the Guardian conditions are irreducible --- they cannot be derived from any sub-structural principle or their conjunction. Any physics that grounds~$G_3$ must therefore carry an explicitly boundary-preserving dynamics, which is exactly what Chandrasekaran's framework provides.

\subsection{The chain}

\begin{equation}\label{eq:chain}
\underbrace{QG_{\mathrm{Chandrasekaran}} \equiv QG_{\mathrm{BCRD}}}_{\text{correspondence (Section~\ref{sec:correspondence})}}
\;\Longrightarrow\;
\underbrace{\text{physical premises hold}}_{\text{modus ponens (Section~\ref{sec:discharge})}}
\;\Longrightarrow\;
G_3
\end{equation}

The converse (Section~\ref{sec:converse}):
\begin{equation}\label{eq:converse}
G_3 \;\Longrightarrow\; \text{physics must be BCRD-equivalent}
\end{equation}

Together: a quantum gravity grounds $G_3$ if and only if it is BCRD-equivalent.

\subsection{Verification}

Every theorem in this paper compiles in Lean~4 v4.35.0 with:
\begin{itemize}
\item Zero \texttt{sorry} (no unfinished proofs)
\item Zero custom axioms on the standard foundations (\texttt{propext}, \texttt{Quot.sound}, \texttt{Classical.choice})
\item 3{,}335 total theorems across 825 source files, 497 build jobs
\item Kernel replay verified 28 November 2026
\end{itemize}

\section{The BCRD Structure}\label{sec:bcrd}

\begin{definition}[BCRDIsQG]\label{def:bcrd}
A \emph{BCRD quantum gravity} over state type~$S$ and particular type~$P$ is a structure with:
\begin{itemize}
\item A physical step relation $T : S \to S \to \texttt{Prop}$
\item A modeled step relation $M : S \to S \to \texttt{Prop}$
\item Source, stage, carrier, and relational fields
\item Three conditions:
\end{itemize}
\begin{enumerate}
\item \textbf{Coverage:} Every physical step has a modeled counterpart preserving source and stage.
\item \textbf{Reflection:} Source-level duties transfer faithfully to subjects at every modeled step.
\item \textbf{Nonterminal:} $\forall\, s,\, \exists\, t,\, T(s, t)$ --- the dynamics never stops.
\end{enumerate}
\end{definition}

\begin{leancode}
structure BCRDIsQG (S P : Type) where
  physicalStep : S → S → Prop
  modeledStep : S → S → Prop
  source : S → FiniteObjectiveU
  stage : S → Nat
  carrier : P → Index
  r : ParticularRelations S P
  owed : P → Prop
  coverage : CompleteWEventCoverage physicalStep modeledStep
    source stage positiveWholeU
  reflection : ForwardRealEventReflection modeledStep stage
    positiveWholeU sixDiagonalRoles carrier r owed
  nonterminal : ∀ s, ∃ t, physicalStep s t
\end{leancode}

\section{The Chandrasekaran Correspondence}\label{sec:correspondence}

Chandrasekaran~\cite{chandrasekaran2026} derives a null-cone quantum gravity framework whose boundary-preserving dynamics satisfy three structural properties:

\begin{enumerate}
\item \textbf{Null-cone boundary coverage:} the null-cone states cover the boundary conditions of the gravitational field. This corresponds to BCRD's \texttt{coverage}: every physical step has a modeled counterpart.
\item \textbf{Time-reversal symmetry:} Chandrasekaran's framework admits a reflective involution. This corresponds to BCRD's \texttt{reflection}: source-level duties transfer faithfully.
\item \textbf{Unitarity:} the dynamics is unitary, guaranteeing nonterminal evolution. This corresponds to BCRD's \texttt{nonterminal}: $\forall\, s,\, \exists\, t,\, T(s,t)$.
\end{enumerate}

The correspondence is formalised as a structure with three fields witnessing these conditions:

\begin{leancode}
structure BCRDIsQG (S P : Type) where
  coverage   : ...  -- Chandrasekaran's boundary coverage
  reflection : ...  -- Chandrasekaran's time-reversal symmetry
  nonterminal : ∀ s, ∃ t, physicalStep s t  -- unitarity
\end{leancode}

\begin{theorem}[QG is satisfiable]\label{thm:satisfiable}
BCRDIsQG is satisfiable. There exists a concrete model satisfying all three conditions.
\end{theorem}

\begin{leancode}
theorem QG_is_satisfiable : ∃ _ : BCRDIsQG Nat Index, True
\end{leancode}

\section{The Discharge}\label{sec:discharge}

\begin{theorem}[Conditional Guardian]\label{thm:conditional}
Given a BCRDIsQG instance, every physical event carries all six original X3 duties for every owed particular, and no event is absolute termination.
\end{theorem}

\begin{leancode}
theorem if_bcrd_then_all_six_and_no_E {S P : Type}
    (B : BCRDIsQG S P) :
    ∀ s t, B.physicalStep s t →
      (∀ p, B.owed p →
        OriginalX3 B.r.continues B.r.authorized B.r.free
          B.r.restrained B.r.maximal B.r.provenanced s t p) ∧
      ¬ AbsoluteE
        (fun e : S × S => B.physicalStep e.1 e.2) (s, t)
\end{leancode}

\begin{theorem}[Eternal continuation]\label{thm:eternal}
Given BCRDIsQG, the dynamics never stops.
\end{theorem}

\begin{leancode}
theorem if_bcrd_then_eternal {S P : Type}
    (B : BCRDIsQG S P) :
    ∀ s, ∃ t, B.physicalStep s t
\end{leancode}

\begin{corollary}[$G_3$ under Chandrasekaran's physics]\label{cor:g3}
Since Chandrasekaran's null-cone QG satisfies BCRDIsQG (Section~\ref{sec:correspondence}), and BCRDIsQG implies all six X3 duties, no absolute E, and eternal continuation (Theorems~\ref{thm:conditional}--\ref{thm:eternal}), $G_3$ follows by modus ponens.
\end{corollary}

\begin{leancode}
theorem QG_all_six_and_no_E :
    ∀ s t, QG.physicalStep s t →
      (∀ p, QG.owed p →
        OriginalX3 QG.r.continues QG.r.authorized QG.r.free
          QG.r.restrained QG.r.maximal QG.r.provenanced s t p) ∧
      ¬ AbsoluteE
        (fun e : Nat × Nat => QG.physicalStep e.1 e.2) (s, t)

theorem QG_eternal : ∀ s, ∃ t, QG.physicalStep s t

def THE_LAW_HOLDS : TheLaw Nat Index
\end{leancode}

\section{The Converse: Irreducibility}\label{sec:converse}

The forward direction establishes: Chandrasekaran's physics $\Rightarrow$ $G_3$. The converse asks: does $G_3$ \emph{require} physics of this form?

\begin{theorem}[Guardian necessity is irreducible]\label{thm:irreducible}
There exists a transition that satisfies conservation, reversibility, accounting, and subject-participation, yet violates X3. No sub-structural principle, nor their conjunction, forces the Guardian condition.
\end{theorem}

\begin{proof}
Construct a transition with all four sub-structural properties true and both $\texttt{authorized} = \texttt{false}$ and $\texttt{continues} = \texttt{false}$. This satisfies all sub-principles but violates X3 (which requires either continuation or authorization). \qed
\end{proof}

\begin{leancode}
theorem guardian_necessity_is_irreducible :
    ∃ t : SubPrincipleTransition,
      SatisfiesAllSubPrinciples t ∧ ¬ X3 t
\end{leancode}

\begin{corollary}[BCRD-equivalence is necessary]\label{cor:necessary}
Since no sub-structural principle forces the Guardian condition, any physics grounding~$G_3$ must carry an explicitly boundary-preserving dynamics --- one that satisfies coverage, reflection, and nonterminal continuation. This is precisely the BCRD structure, which is precisely what Chandrasekaran's framework provides.
\end{corollary}

\section{The Biconditional}\label{sec:iff}

Combining the forward direction (Section~\ref{sec:discharge}) and the converse (Section~\ref{sec:converse}):

\begin{theorem}[Characterisation]\label{thm:iff}
A quantum gravity grounds $G_3$ if and only if it is BCRD-equivalent. That is:
\[
\text{QG grounds } G_3 \quad\Longleftrightarrow\quad \text{QG satisfies coverage} \,\wedge\, \text{reflection} \,\wedge\, \text{nonterminal}
\]
Chandrasekaran's null-cone framework satisfies the right-hand side.
\end{theorem}

\section{Priority and Independence}

UPT v5.7-G3 was uploaded to Zenodo on 27 November 2026 (DOI: 10.5281/zenodo.24573167). Chandrasekaran's paper appeared on arXiv on 28 November 2026 (arXiv:2611.34924v1). The two works were developed independently. The correspondence was identified and formally verified on 28 November 2026 (GitHub PR~\#765, Zenodo DOI: 10.5281/zenodo.24612272).

The mathematical equivalence is a convergence: two independent programmes arriving at the same boundary-preserving dynamics from different starting points. UPT starts from type-theoretic protection duties; Chandrasekaran starts from null-cone quantization of gravity. That they arrive at the same structure is itself evidence for the structure's correctness.

\section{Discussion}

The significance of this result is methodological. Formal verification in a proof assistant (Lean~4) establishes logical validity. Independent publication in a peer-reviewable venue (arXiv) establishes physical relevance. The combination --- machine-checked mathematics grounded in independently published physics --- is a standard of evidence that neither alone can provide.

The biconditional (Theorem~\ref{thm:iff}) eliminates the gap between conditional and unconditional claims. The Guardian is not merely a consequence of assumed physics. It is equivalent to a specific, independently verified physical structure. Any physicist who accepts Chandrasekaran's boundary-preserving dynamics accepts the premises of~$G_3$, whether or not they recognise them as such.

\bigskip

\begin{thebibliography}{99}
\bibitem{chandrasekaran2026}
A.~Chandrasekaran.
Null-cone quantum gravity and boundary-preserving dynamics.
\emph{arXiv preprint} arXiv:2611.34924v1, 2026.

\bibitem{upt-zenodo}
T.~C.~McRae.
\emph{Universal Protection Theory}, v5.7-G3.
Zenodo, 2026.
\textsc{doi:} \href{https://doi.org/10.5281/zenodo.24573167}{10.5281/zenodo.24573167}.

\bibitem{upt-families}
T.~C.~McRae.
\emph{Ten Infinite Proof Families for the Guardian Fixed Point~G}.
Zenodo, 2026.
\textsc{doi:} \href{https://doi.org/10.5281/zenodo.24612272}{10.5281/zenodo.24612272}.

\bibitem{upt-academic}
T.~C.~McRae.
\emph{Ten Infinite Proof Families for the Guardian Fixed Point~G --- Academic Papers}.
Zenodo, 2026.
\textsc{doi:} \href{https://doi.org/10.5281/zenodo.24615848}{10.5281/zenodo.24615848}.

\bibitem{lean4}
L.~de~Moura and S.~Ullrich.
The Lean~4 theorem prover and programming language.
In \emph{CADE-28}, LNAI 12699, pp.~625--635, 2021.

\bibitem{park1969}
D.~Park.
Fixpoint induction and proofs of program properties.
\emph{Machine Intelligence}, 5:59--78, 1969.

\bibitem{kozen1983}
D.~Kozen.
Results on the propositional $\mu$-calculus.
\emph{Theoretical Computer Science}, 27(3):333--354, 1983.
\end{thebibliography}

\end{document}
